\section{Pairwise contrast audit}
\label{suppsec:contrasts}

\subsection{Recomputation of the contrasts}

The matched contrasts are computed directly from the absolute scalar values:
\begin{equation}
I_D(q_i,q_j) = \frac{D^*(q_j)-D^*(q_i)}{D^*(q_i)}\times 100\%, \qquad
I_P(q_i,q_j) = \frac{P^*(q_j)-P^*(q_i)}{P^*(q_i)}\times 100\%.
\end{equation}
The sign convention is fixed: a positive $I_D$ indicates that the higher-exponent geometry dissipates more energy than the lower-exponent one, relative to the latter.

\subsection{Sign and normalization conventions}

The denominator in the percentage is always $D^*(q_i)$ (or $P^*(q_i)$) of the earlier configuration in the pair. The three transitions studied are $(2,4)$, $(4,6)$, and $(6,8)$; the full sequence is available at all three Reynolds numbers. The full contrast dataset for all resolutions is provided in Table~\ref{supptab:contrasts}, which also lists, for every transition and amplitude, the span of $I_D$ over $N=128$, $256$, $384$, and $512$ (the last resolution available only at $Re=400$ and $Re=1000$, $\alpha=0$). That span is the resolution floor of the contrast itself, and the working criterion throughout is that a contrast is resolved only when $|I_D|$ at the finest available resolution exceeds it.

At $Re=400$ the $(6,8)$ contrast at $\alpha=0$ is $+6\times 10^{-5}\%$ against a floor of $1.1\times 10^{-4}$ percentage points, so it is reported only as an upper bound and its sign is not stable across grids. The same transition at $\alpha=0.1$ ($-0.0017\%$, floor $8.5\times 10^{-5}$) and the $(4,6)$ contrast at both amplitudes ($+0.0129\%$, floor $5.0\times 10^{-4}$, and $-0.0039\%$, floor $8.7\times 10^{-4}$) are resolved, the spans being $5\%$, $3.9\%$ and $23\%$ of their magnitudes. At $Re=1000$ both trailing transitions are resolved: $(4,6)=+0.0103\%$ and $(6,8)=-0.0037\%$ at $\alpha=0$, with spans over $N=128$--$512$ of $6.9\times10^{-4}$ and $1.1\times10^{-4}$ percentage points, and $(4,6)=-0.0242\%$, $(6,8)=-0.0052\%$ at $\alpha=0.1$, with spans of $2.2\times10^{-3}$ and $5.5\times10^{-4}$ percentage points. At the undulated wall the two trailing transitions are already negative at $Re=400$ ($-0.0039\%$ and $-0.0017\%$), each larger in magnitude than its own span. The $N=512$ data at $Re=1000$ confirm the monotonicity of the $(2,4)$ contrast ($+0.592\%$) and strengthen the negative sign of the $(6,8)$ transition ($-0.0036\%$).

At $Re=100$ the trailing transitions are $+0.00296\%$ and $+0.00042\%$ at $\alpha=0$, and $+0.00098\%$ and $+0.00006\%$ at $\alpha=0.1$. Only part of this set is directly usable. On the two coarser grids the $n\ge 6$ cases at this Reynolds number run at $Ma_{\max}=0.18$ and $0.09$ while their $n=4$ partners run at $0.068$ and $0.034$, so the coarse members of the $(4,6)$ sequence are not in the same compressibility regime as the reference; the $1.145$ percentage-point apparent span therefore measures that mismatch rather than grid convergence, and no resolution floor can be assigned to the $Re=100$ $(4,6)$ contrast, whose $N=384$ value is quoted as an upper bound. The $(6,8)$ pairs match in Mach number on every grid, their compressibility error cancels in common mode, and their spans ($8.1\times 10^{-5}$ and $5.5\times 10^{-5}$ percentage points) are genuine floors: $+0.00042\%$ at $\alpha=0$ is resolved at $5.2$ times its span, while $+0.00006\%$ at $\alpha=0.1$ is marginal at $1.1$ times its span and is quoted for completeness only.

\begin{table}[!htbp]
\footnotesize
\centering
\renewcommand{\arraystretch}{0.92}
\setlength{\tabcolsep}{3.5pt}
\caption{Full contrast dataset for all resolutions, transitions and amplitudes. $\delta_{IP/ID}=100|I_P/I_D-1|$ is the relative estimator discrepancy and the second-to-last column is the absolute power-minus-dissipation contrast difference, both in percentage points (pp). The last column, printed once per transition and amplitude on its $N=384$ row, is the \emph{own resolution span} of that contrast: the range of $I_D$ over $N=128$, $256$, $384$, and $512$ (the last resolution available only at $Re=400$ and $Re=1000$, $\alpha=0$). A contrast is treated as resolution-resolved only when $|I_D|$ at the finest available resolution exceeds its own span; by that criterion every $(2,4)$ contrast at $\alpha=0$, the trailing $(4,6)$ and $(6,8)$ contrasts at $Re=1000$ for both amplitudes, $(4,6)$ at $Re=400$ and $\alpha=0$, and $(6,8)$ at $Re=400$ and $\alpha=0.1$ are resolved, whereas $Re=400$, $(6,8)$, $\alpha=0$ ($0.00006\%$ against a $0.00011$~pp span) is not and is quoted only as an upper bound. The $Re=100$ trailing spans require a separate caveat: at $N=128$ and $256$ the $n\ge6$ cases sit at local Mach number $Ma_{\max}\simeq 0.18$ and $0.09$ while the $n=4$ reference is at $0.068$ and $0.034$, so the $(4,6)$ pairs compare different compressibility regimes and their spans ($1.145$~pp) measure that mismatch rather than grid convergence. The $(6,8)$ pairs match in Mach number at every $N$, their errors cancel in common mode, and their spans ($8\times10^{-5}$ and $5\times10^{-5}$~pp) are genuine resolution floors; the $Re=100$, $(4,6)$ values at $N=384$ are accordingly quoted as upper bounds. Entries with $|I_D|<10^{-4}\%$ have a numerically meaningless $\delta_{IP/ID}$ and are marked ``--''.}
\label{supptab:contrasts}
\begin{tabular}{@{}ccccccccc@{}}
\hline
$Re$ & $N$ & Transition & $\alpha$ & $I_D$ [\%] & $I_P$ [\%] & $\delta_{IP/ID}$ [\%] & $I_P-I_D$ [pp] & span [pp] \\
\hline
100  & 128 & $(2,4)$ & 0.00 & $+0.02161$ & $+0.01971$ &  8.77 & $-0.00189$ &  \\
100  & 256 & $(2,4)$ & 0.00 & $+0.02195$ & $+0.02079$ &  5.27 & $-0.00116$ &  \\
100  & 384 & $(2,4)$ & 0.00 & $+0.02178$ & $+0.02110$ &  3.14 & $-0.00068$ & $+0.00034$ \\
\hline
100  & 128 & $(2,4)$ & 0.10 & $+0.01078$ & $+0.00587$ & 45.60 & $-0.00492$ &  \\
100  & 256 & $(2,4)$ & 0.10 & $+0.01097$ & $+0.00675$ & 38.42 & $-0.00421$ &  \\
100  & 384 & $(2,4)$ & 0.10 & $+0.00806$ & $+0.00705$ & 12.63 & $-0.00102$ & $+0.00290$ \\
\hline
100  & 128 & $(4,6)$ & 0.00 & $+1.14810$ & $+1.26723$ & 10.38 & $+0.11913$ &  \\
100  & 256 & $(4,6)$ & 0.00 & $+0.50079$ & $+0.81674$ & 63.09 & $+0.31595$ &  \\
100  & 384 & $(4,6)$ & 0.00 & $+0.00296$ & $+0.00264$ & 10.63 & $-0.00031$ & $+1.14514$ \\
\hline
100  & 128 & $(4,6)$ & 0.10 & $+1.14661$ & $+1.26424$ & 10.26 & $+0.11763$ &  \\
100  & 256 & $(4,6)$ & 0.10 & $+0.49890$ & $+0.81430$ & 63.22 & $+0.31540$ &  \\
100  & 384 & $(4,6)$ & 0.10 & $+0.00098$ & $+0.00073$ & 25.60 & $-0.00025$ & $+1.14563$ \\
\hline
100  & 128 & $(6,8)$ & 0.00 & $+0.00034$ & $+0.00036$ &  5.74 & $+0.00002$ &  \\
100  & 256 & $(6,8)$ & 0.00 & $+0.00040$ & $+0.00039$ &  4.43 & $-0.00002$ &  \\
100  & 384 & $(6,8)$ & 0.00 & $+0.00042$ & $+0.00040$ &  4.46 & $-0.00002$ & $+0.00008$ \\
\hline
100  & 128 & $(6,8)$ & 0.10 & $+0.00001$ & $+0.00003$ & -- & $+0.00003$ &  \\
100  & 256 & $(6,8)$ & 0.10 & $+0.00003$ & $+0.00004$ & -- & $+0.00000$ &  \\
100  & 384 & $(6,8)$ & 0.10 & $+0.00006$ & $+0.00005$ & -- & $-0.00001$ & $+0.00005$ \\
\hline
400  & 128 & $(2,4)$ & 0.00 & $+0.32251$ & $+0.33406$ &  3.58 & $+0.01155$ &  \\
400  & 256 & $(2,4)$ & 0.00 & $+0.32462$ & $+0.32990$ &  1.63 & $+0.00528$ &  \\
400  & 384 & $(2,4)$ & 0.00 & $+0.32488$ & $+0.32825$ &  1.04 & $+0.00337$ &  \\
400  & 512 & $(2,4)$ & 0.00 & $+0.32485$ & $+0.32711$ &  0.70 & $+0.00226$ & $+0.00237$ \\
\hline
400  & 128 & $(2,4)$ & 0.10 & $+0.14269$ & $+0.13580$ &  4.83 & $-0.00689$ &  \\
400  & 256 & $(2,4)$ & 0.10 & $+0.14518$ & $+0.13404$ &  7.68 & $-0.01114$ &  \\
400  & 384 & $(2,4)$ & 0.10 & $+0.13326$ & $+0.13352$ &  0.20 & $+0.00026$ & $+0.01192$ \\
\hline
400  & 128 & $(4,6)$ & 0.00 & $+0.01237$ & $+0.01117$ &  9.71 & $-0.00120$ &  \\
400  & 256 & $(4,6)$ & 0.00 & $+0.01287$ & $+0.01142$ & 11.28 & $-0.00145$ &  \\
400  & 384 & $(4,6)$ & 0.00 & $+0.01287$ & $+0.01155$ & 10.28 & $-0.00132$ & $+0.00050$ \\
\hline
400  & 128 & $(4,6)$ & 0.10 & $-0.00474$ & $-0.00465$ &  1.73 & $+0.00008$ &  \\
400  & 256 & $(4,6)$ & 0.10 & $-0.00456$ & $-0.00444$ &  2.68 & $+0.00012$ &  \\
400  & 384 & $(4,6)$ & 0.10 & $-0.00386$ & $-0.00432$ & 11.85 & $-0.00046$ & $+0.00087$ \\
\hline
400  & 128 & $(6,8)$ & 0.00 & $-0.00005$ & $-0.00037$ & -- & $-0.00032$ &  \\
400  & 256 & $(6,8)$ & 0.00 & $+0.00005$ & $-0.00023$ & -- & $-0.00029$ &  \\
400  & 384 & $(6,8)$ & 0.00 & $+0.00006$ & $-0.00015$ & -- & $-0.00021$ & $+0.00011$ \\
\hline
400  & 128 & $(6,8)$ & 0.10 & $-0.00160$ & $-0.00170$ &  6.64 & $-0.00011$ &  \\
400  & 256 & $(6,8)$ & 0.10 & $-0.00163$ & $-0.00174$ &  6.46 & $-0.00011$ &  \\
400  & 384 & $(6,8)$ & 0.10 & $-0.00168$ & $-0.00176$ &  4.90 & $-0.00008$ & $+0.00008$ \\
\hline
1000 & 128 & $(2,4)$ & 0.00 & $+0.57758$ & $+0.61609$ &  6.67 & $+0.03851$ &  \\
1000 & 256 & $(2,4)$ & 0.00 & $+0.58769$ & $+0.60521$ &  2.98 & $+0.01752$ &  \\
1000 & 384 & $(2,4)$ & 0.00 & $+0.59087$ & $+0.60217$ &  1.91 & $+0.01130$ &  \\
1000 & 512 & $(2,4)$ & 0.00 & $+0.59178$ & $+0.59647$ &  0.79 & $+0.00469$ & $+0.01418$ \\
\hline
1000 & 128 & $(2,4)$ & 0.10 & $+0.10893$ & $+0.08443$ & 22.49 & $-0.02449$ &  \\
1000 & 256 & $(2,4)$ & 0.10 & $+0.11989$ & $+0.08010$ & 33.20 & $-0.03980$ &  \\
1000 & 384 & $(2,4)$ & 0.10 & $+0.07824$ & $+0.07971$ &  1.87 & $+0.00147$ & $+0.04165$ \\
\hline
1000 & 128 & $(4,6)$ & 0.00 & $+0.01074$ & $+0.00532$ & 50.44 & $-0.00542$ &  \\
1000 & 256 & $(4,6)$ & 0.00 & $+0.01058$ & $+0.00577$ & 45.42 & $-0.00481$ &  \\
1000 & 384 & $(4,6)$ & 0.00 & $+0.01033$ & $+0.00615$ & 40.41 & $-0.00417$ &  \\
1000 & 512 & $(4,6)$ & 0.00 & $+0.01005$ & $+0.00647$ & 35.62 & $-0.00358$ & $+0.00069$ \\
\hline
1000 & 128 & $(4,6)$ & 0.10 & $-0.02589$ & $-0.02645$ &  2.19 & $-0.00057$ &  \\
1000 & 256 & $(4,6)$ & 0.10 & $-0.02643$ & $-0.02571$ &  2.72 & $+0.00072$ &  \\
1000 & 384 & $(4,6)$ & 0.10 & $-0.02424$ & $-0.02537$ &  4.70 & $-0.00114$ & $+0.00220$ \\
\hline
1000 & 128 & $(6,8)$ & 0.00 & $-0.00371$ & $-0.00460$ & 23.83 & $-0.00088$ &  \\
1000 & 256 & $(6,8)$ & 0.00 & $-0.00360$ & $-0.00443$ & 22.97 & $-0.00083$ &  \\
1000 & 384 & $(6,8)$ & 0.00 & $-0.00371$ & $-0.00429$ & 15.60 & $-0.00058$ &  \\
1000 & 512 & $(6,8)$ & 0.00 & $-0.00364$ & $-0.00419$ & 15.11 & $-0.00055$ & $+0.00011$ \\
\hline
1000 & 128 & $(6,8)$ & 0.10 & $-0.00482$ & $-0.00536$ & 11.19 & $-0.00054$ &  \\
1000 & 256 & $(6,8)$ & 0.10 & $-0.00461$ & $-0.00535$ & 16.18 & $-0.00075$ &  \\
1000 & 384 & $(6,8)$ & 0.10 & $-0.00516$ & $-0.00557$ &  8.06 & $-0.00042$ & $+0.00055$ \\
\hline
\end{tabular}
\end{table}


